%%%%%%%%%%%%%%%%%
% This is an sample CV template created using altacv.cls
% (v1.7.4, 30 July 2025) written by LianTze Lim (liantze@gmail.com). Compiles with pdfLaTeX, XeLaTeX and LuaLaTeX.
%
%% It may be distributed and/or modified under the
%% conditions of the LaTeX Project Public License, either version 1.3
%% of this license or (at your option) any later version.
%% The latest version of this license is in
%%    http://www.latex-project.org/lppl.txt
%% and version 1.3 or later is part of all distributions of LaTeX
%% version 2003/12/01 or later.
%%%%%%%%%%%%%%%%

%% Use the "normalphoto" option if you want a normal photo instead of cropped to a circle
% \documentclass[10pt,a4paper,withhyper,normalphoto]{altacv}

\documentclass[10pt,a4paper,withhyper]{altacv}
%% AltaCV uses the fontawesome5 and simpleicons packages.
%% See http://texdoc.net/pkg/fontawesome5 and http://texdoc.net/pkg/simpleicons for full list of symbols.

% Change the page layout if you need to
\geometry{left=1.25cm,right=1.25cm,top=1.5cm,bottom=1.5cm,columnsep=1.2cm}

% The paracol package lets you typeset columns of text in parallel
\usepackage{paracol}

% Change the font if you want to, depending on whether
% you're using pdflatex or xelatex/lualatex
% WHEN COMPILING WITH XELATEX PLEASE USE
% xelatex -shell-escape -output-driver="xdvipdfmx -z 0" sample.tex
\iftutex
  % If using xelatex or lualatex:
  \setmainfont{Roboto Slab}
  \setsansfont{Lato}
  \renewcommand{\familydefault}{\sfdefault}
\else
  % If using pdflatex:
  \usepackage[rm]{roboto}
  \usepackage[defaultsans]{lato}
  % \usepackage{sourcesanspro}
  \renewcommand{\familydefault}{\sfdefault}
\fi

% Change the colours if you want to
\definecolor{SlateGrey}{HTML}{2E2E2E}
\definecolor{LightGrey}{HTML}{666666}
\definecolor{PastelRed}{HTML}{8F0D0D}
\definecolor{MainColor}{HTML}{28555e}
\definecolor{LightMainColor}{HTML}{3a6b75}
\definecolor{GoldenEarth}{HTML}{E7D192}
\definecolor{DarkGoldenEarth}{HTML}{d18c5c}
\colorlet{name}{black}
\colorlet{tagline}{MainColor}
\colorlet{heading}{MainColor}
\colorlet{headingrule}{GoldenEarth}
\colorlet{subheading}{DarkGoldenEarth}
\colorlet{accent}{DarkGoldenEarth}
\colorlet{emphasis}{SlateGrey}
\colorlet{body}{LightGrey}

% Change some fonts, if necessary
\renewcommand{\namefont}{\Huge\rmfamily\bfseries}
\renewcommand{\personalinfofont}{\footnotesize}
\renewcommand{\cvsectionfont}{\LARGE\rmfamily\bfseries}
\renewcommand{\cvsubsectionfont}{\large\bfseries}


% Change the bullets for itemize and rating marker
% for \cvskill if you want to
\renewcommand{\cvItemMarker}{{\small\textbullet}}
\renewcommand{\cvRatingMarker}{\faCircle}
% ...and the markers for the date/location for \cvevent
% \renewcommand{\cvDateMarker}{\faCalendar*[regular]}
% \renewcommand{\cvLocationMarker}{\faMapMarker*}


% If your CV/résumé is in a language other than English,
% then you probably want to change these so that when you
% copy-paste from the PDF or run pdftotext, the location
% and date marker icons for \cvevent will paste as correct
% translations. For example Spanish:
% \renewcommand{\locationname}{Ubicación}
% \renewcommand{\datename}{Fecha}


%% Use (and optionally edit if necessary) this .tex if you
%% want to use an author-year reference style like APA(6)
%% for your publication list
% \input{pubs-authoryear}

%% Use (and optionally edit if necessary) this .tex if you
%% want an originally numerical reference style like IEEE
%% for your publication list
\input{pubs-num}

%% sample.bib contains your publications
\addbibresource{sample.bib}

\begin{document}
\name{Yann HALLOUARD}
\tagline{Senior Data Scientist \& ML Engineer}
%% You can add multiple photos on the left or right
\photoR{2.8cm}{profile_picture}
% \photoL{2.5cm}{Yacht_High,Suitcase_High}

\personalinfo{%
  % Not all of these are required!
  \email{yann.hal@gmail.com}
  \phone{+33 6 35 43 14 53}
  \mailaddress{27 av Villemain, 75014, Paris}
  \location{Paris, France}
  % \homepage{www.homepage.com}
  % \twitter{@twitterhandle}
  % \xtwitter{@x-handle}
  \linkedin{in/yann-hallouard}
  \github{YHallouard}
  \printinfo{\faMedium}{@yann-hal}[https://medium.com/@yann-hal]
  % \orcid{0000-0000-0000-0000}
  %% You can add your own arbitrary detail with
  %% \printinfo{symbol}{detail}[optional hyperlink prefix]
  % \printinfo{\faPaw}{Hey ho!}[https://example.com/]

  %% Or you can declare your own field with
  %% \NewInfoFiled{fieldname}{symbol}[optional hyperlink prefix] and use it:
  % \NewInfoField{gitlab}{\faGitlab}[https://gitlab.com/]
  % \gitlab{your_id}
  %%
  %% For services and platforms like Mastodon where there isn't a
  %% straightforward relation between the user ID/nickname and the hyperlink,
  %% you can use \printinfo directly e.g.
  % \printinfo{\faMastodon}{@username@instace}[https://instance.url/@username]
  %% But if you absolutely want to create new dedicated info fields for
  %% such platforms, then use \NewInfoField* with a star:
  % \NewInfoField*{mastodon}{\faMastodon}
  %% then you can use \mastodon, with TWO arguments where the 2nd argument is
  %% the full hyperlink.
  % \mastodon{@username@instance}{https://instance.url/@username}
}

\makecvheader
%% Depending on your tastes, you may want to make fonts of itemize environments slightly smaller
% \AtBeginEnvironment{itemize}{\small}

%% Set the left/right column width ratio to 6:4.
\columnratio{0.6}

% Start a 2-column paracol. Both the left and right columns will automatically
% break across pages if things get too long.
\begin{paracol}{2}
\cvsection{Experience}

\cvevent{Senior Data Scientist}{TotalEnergies Digital Factory}{Apr 2025 -- Present}{Paris}
\textbf{DataTeam - InnovationLab: Co-produced AI products with Mistral}
\begin{itemize}
\item LLM production pipelines: Robust generative AI industrialization
\item Parallel implementation: Vocal chat service for end users \& report generator for chemical refinery
\item Technology stack: Temporal.io, Kubernetes, Entity Matching, Production Pipelines
\end{itemize}

\divider

\cvevent{Senior Data Scientist}{TotalEnergies Digital Factory}{Feb 2023 -- Apr 2025}{Paris}
\textbf{SAFEWORLD - Real-time Industrial Safety System}
\begin{itemize}
\item AI-powered safety compliance monitoring for 10 industrial sites
\item Computer Vision optimization: RT-DETR, DinoV2 for real-time inference
\item Edge computing architecture: 8-second event-to-platform detection latency
\item Model performance: mAP@0.5 of 0.8 (35m), 0.6 (70m)
\item Technology stack: DynamoDB, AWS Lambda, Serverless Architecture
\end{itemize}

\divider

\cvevent{Data Scientist}{TotalEnergies Digital Factory}{Nov 2021 -- Jan 2023}{Paris}
\textbf{DEGAS - Topographic Signal Denoising Platform}
\begin{itemize}
\item ML denoising system for landslide risk monitoring sensors
\item Change Point Detection: 96\% accuracy on multivariate time series
\item Autoencoder-based correction with continuous deployment
\item Technology stack: Azure ML, TimescaleDB, Azure Functions
\end{itemize}

\divider

\cvevent{Data Scientist}{TotalEnergies Digital Factory}{Apr 2021 -- Nov 2021}{Paris}
\textbf{RADAR - Oil Well Production Anomaly Detection}
\begin{itemize}
\item Bayesian anomaly detection for production forecasting enhancement
\item Operational noise filtering from genuine production issues
\item Technology stack: Databricks, Azure Data Factory, MLflow, PostgreSQL
\end{itemize}

\divider

\cvevent{Junior Data Scientist}{TotalEnergies Digital Factory}{Jun 2020 -- Apr 2021}{Paris}
\begin{itemize}
\item Energy efficiency analysis via electricity consumption clustering
\item FPSO CO2 optimization using causal analysis and multivariate classification
\item Helicopter crew rotation optimization (VRP with time windows)
\item Technology stack: Python, Streamlit, OR-Tools, Statistical Modeling
\end{itemize}

\medskip

\cvsection{Projects}

\cvevent{Héméa - Mobile Health App}{Personal Project}{May 2025 -- Present}{}
\begin{itemize}
\item Mobile application for organizing, storing and visualizing biological analyses
\item User-friendly interface for medical test results management
\item Data visualization dashboard for health monitoring trends
\item Secure storage and privacy-focused architecture
\item Technology stack: Mistral, React Native, Expo Go, SQLite
\end{itemize}

%% \divider
%% 
%% \cvevent{Project 2}{Funding agency/institution}{Project duration}{}
%% A short abstract would also work.

\medskip

\cvsection{Publications}

%% Specify your last name(s) and first name(s) as given in the .bib to automatically bold your own name in the publications list.
%% One caveat: You need to write \bibnamedelima where there's a space in your name for this to work properly; or write \bibnamedelimi if you use initials in the .bib
%% You can specify multiple names, especially if you have changed your name or if you need to highlight multiple authors.
\mynames{Yann\bibnamedelima Hallouard}
%% MAKE SURE THERE IS NO SPACE AFTER THE FINAL NAME IN YOUR \mynames LIST

\nocite{*}

% \printbibliography[heading=pubtype,title={\printinfo{\faBook}{Books}},type=book]
% 
% \divider

\printbibliography[heading=pubtype,title={\printinfo{\faFile*[regular]}{Articles}},type=article]

% \divider
% 
% \printbibliography[heading=pubtype,title={\printinfo{\faUsers}{Conference Proceedings}},type=inproceedings]

%% Switch to the right column. This will now automatically move to the second
%% page if the content is too long.
\switchcolumn

\cvsection{My Life Philosophy}

\begin{quote}
``Graduated from École des Mines de Saint-Étienne, passionate about transforming models into tangible and usable solutions. With solid experience in production deployment, motivated to continue developing products that bring real added value to users.''
\end{quote}

\cvsection{Most Proud of}

\cvachievement{\faTrophy}{TDF GenAI Hackathon}{2nd Place - Oct 2024}

\divider

\cvachievement{\faCertificate}{Certifications}{Azure - DP200 - Sept 2022 \newline Hugging Face - Fondamental of MCP - June 2025}

\divider

\cvachievement{\faUsers}{Assistant Secretary CSE}{TotalEnergies Digital Factory - Mar 2022 to present - Managing budget and employee programs for 150+ Digital Factory employees}

\divider

\cvachievement{\faCode}{Community Contributions}{Sktime, Moto, Roboflow}

\medskip

\cvsection{Mentoring}

\begin{itemize}
\item \textbf{Team \& intern supervision}: Data technical leadership on delevered products and intern mentoring across multiple projects
\item \textbf{Technical community animation}: Organizing community of practices within Digital Factory
\item \textbf{Recruitment \& interviews}: Conducted 30+ technical interviews for Data Science and ML Engineering positions
\item \textbf{Training \& seminars}: Delivered technical presentations to audiences up to 300 people
\end{itemize}


\cvsection{Skills}

{\small
\cvtag{Community Building}
\cvtag{Product}
\cvtag{Deep Learning}
\cvtag{Computer Vision}
\cvtag{LLMs}
\cvtag{MLOps}
\cvtag{Statistics}
\cvtag{Transformers}
\cvtag{Python}
\cvtag{AWS}
\cvtag{Azure}
\cvtag{Temporal.io}
\cvtag{k9s}
\cvtag{Clean Code}
\cvtag{Edge Computing}
\cvtag{Serverless}
\cvtag{Real-time Systems}
\cvtag{User centric}
}

\medskip
\newpage

\cvsection{Languages}

\cvskill{Français}{5}
\divider

\cvskill{English - TOEIC 825}{4.5}

%% Yeah I didn't spend too much time making all the
%% spacing consistent... sorry. Use \smallskip, \medskip,
%% \bigskip, \vspace etc to make adjustments.
\medskip

\cvsection{Education}

\cvevent{Engineering Degree}{École Nationale Supérieure des Mines de Saint-Étienne}{Completed in 2019}{}
Data Science, Big Data, ROAD specialization\\
Gestion et Finance

\divider

\cvevent{Master of Business Intelligence and Innovation Management}{IAE Saint-Étienne}{Completed in 2019}{}

\divider

\cvevent{ML in Embedded Systems}{Seoul National University}{Completed in 2019}{}

\divider

\cvevent{Classe Préparatoire - PCSI/PSI}{Lycée Bellevue}{Completed in 2016}{}

\cvsection{Referees}

\cvref{Marc HOBBALLAH}{TDF / Lead Data Scientist}{marc.hobballah@totalenergies.com}
{TotalEnergies Digital Factory\\Paris, France}

\divider

\cvref{Timothée RETAILLEAU}{Accenture / AI \& Data Associate Manager}{ti.retailleau@gmail.com}
{Accenture\\Paris, France}


\end{paracol}


\end{document}
